\section{Two definitions of nearby cycles}

On $\mathbb{A}_{\mathbb{C}}^1$ with coordinate $t$, given any $\alpha \in \mathbb{C}$, we have a left indecomposable $\D_{{A}^1_\C}$ module
\[
    K^\alpha_m :=  \frac{\D_{\mathbb{A}^1_\C}}{\D_{\mathbb{A}^1_\C}\cdot(t\partial_t - \alpha)^m}
\]
First note that $t\in K^\alpha_m$ has \textbf{both} a left and a right inverse, i.e., action by multiplying $t$ induce a bijection.
So $K^\alpha_m$ not only a (left, see 2. in Lem.~\ref{lem:quasi-coherent-sheaf-and-localization-of-D-modules}) $\C[t]\subseteq \D_{\mathbb{A}^1_\C}$ module, but also a (left) $\C[t,t^{-1}]$-module that extending the $\C[t]$-module structure.

We take sections $e^\alpha_i := \overline{(t\partial_t - \alpha)^{m-1-i}}$ (equivalence class) then we have a presentation:
\[
    K^\alpha_m = \bigoplus_{i=0}^{m-1} \mathbb{C}[t, t^{-1}] e^\alpha_i
\]
with $\partial_t e^\alpha_i = \frac{1}{t}(\alpha e^\alpha_i + e^\alpha_{i-1})$ for all $i$ (note we have claimed that $t$ is invertible), where $e^\alpha_{-1} = 0$,
and the direct sum notation should be treated as \textbf{internal} direct sum.

Then when restricted to $\mathbb{A}_{\mathbb{C}}^1 \setminus \{0\}$, $K^\alpha_m\mid_{\mathbb{A}_{\mathbb{C}}^1 \setminus \{0\}}$
becomes a \textbf{finitely generated}, (locally) free $\mathcal{O}_{\mathbb{A}_{\mathbb{C}}^1 \setminus \{0\}}$-module, i.e., an \textbf{integrable connection} (see \hyperref[def:integrable-connection]{definition}).

Then the matrix (note $K^\alpha_m$ is not a vector space since $\C[t,t^{-1}]$ is not a field) of $t\partial_t$-action on $K_m^{\alpha}$ is $J_{\alpha,m}$, the $m\times m$ Jordan block with eigenvalue $\alpha$:
\[
    J_{\alpha,m} = \begin{bmatrix}
        \alpha & 1      &        &        \\
               & \alpha & \ddots &        \\
               &        & \ddots & 1      \\
               &        &        & \alpha
    \end{bmatrix}_{m \times m}
\]

\begin{remark}
    Here $e_\ell^{\alpha}$ can be understood as the formal symbol
    of the multivalued function
    \[
        t^{\alpha} \cdot \frac{(\log t)^\ell}{\ell!}.
    \]
    And $K^\alpha_{m}$ indicates about the section not flat,
    a \textbf{flat section} means the section is in the kernel of connection
    \[
        \nabla:=t\partial_t \in\End_\C(K^\alpha_m).
    \]

    Indeed in \cite[Ch.~5]{hotta2008d},
    they develop arguments or definitions such as ``moderate growth'' by multivalued functions,
    so it is not only as a heuristic symbol to introduce multivalued function.
\end{remark}

Then locally around a simply connected neighborhood (Euclidean topology) of $(\mathbb{A}_{\mathbb{C}}^1 \setminus \{0\})^{\mathrm{an}}$,
with simply connectedness, we obtain that $\mathrm{\log}\, t$ is a well-defined analytic function (so we should consider the analytification of the module $K^\alpha_m$).
Consider
$$
    v_{\ell} := e^{-J_{\alpha,m} \log t} \cdot e_{\ell}^{\alpha} \in (K^\alpha_m)^{\mathrm{an}}
    = \frac{\D^{\mathrm{an}}_{(\mathbb{A}^1_\C)^\mathrm{an}}}{\D^{\mathrm{an}}_{(\mathbb{A}^1_\C)^\mathrm{an}}\cdot(t\partial_t - \alpha)^m}
    = \bigoplus_{i=0}^{m-1} \mathcal{O}^{\mathrm{an}}_{(\mathbb{A}^1_\C)^\mathrm{an}}\left[\frac{1}{t}\right] e^\alpha_i,
$$
for $\ell = 0, \cdots, m-1$,
we then have, by the Leibniz rule,
\begin{align*}
    t\partial_t \circ \left(e^{-J_{\alpha,m} \log t} \cdot e_{\ell}^{\alpha}\right)
     & := t\partial_t \cdot e^{-J_{\alpha,m} \log t} \cdot e_{\ell}^{\alpha}                                                                                                                                                                                             \\
     & = \left[t\partial_t, e^{-J_{\alpha,m} \log t}\right] \cdot e_{\ell}^{\alpha} + e^{-J_{\alpha,m} \log t} \cdot t\partial_t \cdot e_{\ell}^{\alpha}                                                                                                                 \\
     & = t\partial_t\left(e^{-J_{\alpha,m} \log t}\right) \cdot e_{\ell}^{\alpha} + e^{-J_{\alpha,m} \log t} \cdot t\partial_t \cdot e_{\ell}^{\alpha}                              &  & \text{(see \cite[Prop~2.25(iii)]{mustata2023dmodules})}                         \\
     & = t \left( -J_{\alpha,m} \frac{1}{t} e^{-J_{\alpha,m} \log t} \right) \cdot e_{\ell}^{\alpha} + e^{-J_{\alpha,m} \log t} \cdot \left( J_{\alpha,m} e_{\ell}^{\alpha} \right)                                                                                      \\
     & = -J_{\alpha,m} \left( e^{-J_{\alpha,m} \log t} \cdot e_{\ell}^{\alpha} \right) + J_{\alpha,m} \left( e^{-J_{\alpha,m} \log t} \cdot e_{\ell}^{\alpha} \right)               &  & \text{($\C$ is in the \textbf{center} of $\D^{\an}_{(\mathbb{A}^1_\C)^{\an}}$)} \\
     & = 0
\end{align*}
we use $\cdot$ to denote the multiplication in $\D^{\an}_{(\mathbb{A}^1_\C)^{\an}}$, and we get $v_{\ell}$ are $t\partial_t$-flat.

\bigskip

We have a local system, recall the equivalence given in Prop.~\ref{prop:classification-of-local-system}:
$$
    L_m^{\lambda}=\left(\bigoplus \limits_{i=0}^{m-1}\mathbb{C}\cdot v_{i}, e^{-2\pi iJ_{\alpha,m}}\right)
$$
There are several things to be point out:
\begin{enumerate}
    \item Taking fiber (not stalk) at any given point $p\in(\mathbb{A}_{\mathbb{C}}^1 \setminus \{0\})^{\mathrm{an}}$, $\Ker(\nabla)$ becomes a finite ($m$) dimensional $\C$ vector space.
    \item On a simply connected Euclidean open set $p \in U \subset (\mathbb{A}_{\mathbb{C}}^1 \setminus \{0\})^{\mathrm{an}}$, we can give $\log\, t(p)$ a concreate $\C$-value.
    \item View $e^\alpha_i(p)$ as the canonical basis vector (the $i+1$-th entry is $1$, others are $0$)
          \[
              \begin{pmatrix}
                  0 \\ \vdots \\ 1 \\ \vdots \\ 0
              \end{pmatrix}\in\C^m,
          \]
          where the $1$ is at the $(i+1)$-th row. Then we identify
          \[
              v^\alpha_i(p) = e^{-J_{\alpha,m} \log t}(p) \cdot e^\alpha_i(p)
          \]
          with
          \[
              e^{-J_{\alpha,m} \log t}(p) \cdot \begin{pmatrix}
                  0 \\ \vdots \\ 1 \\ \vdots \\ 0
              \end{pmatrix} \in \C^m.
          \]
    \item So the monodromy action (of a generator of $\pi_1((\mathbb{A}_{\mathbb{C}}^1 \setminus \{0\})^{\mathrm{an}}, p)$) acts on $\C^m$ by left multiplication with $e^{-2\pi i J_{\alpha,m}}$.
\end{enumerate}
We denote $\lambda = e^{-2\pi i \alpha}\in \C$.

\bigskip

Since $K^\alpha_m\mid_{\mathbb{A}_{\mathbb{C}}^1 \setminus \{0\}}$ is an integrable connection on $\mathbb{A}_{\mathbb{C}}^1 \setminus \{0\}$ (note that $\C[t, t^{-1}]$ is also smooth finite type over $\C$), recall by Riemann-Hilbert correspondence Thm.~\ref{thm:basic-RH-corr}, we have \cite[Thm.~7.2.1]{hotta2008d} that in $\D^b(\underline{\C}_{X^\mathrm{an}})$,
$$
    \text{DR}\left((K_m^{\alpha}\mid_{\mathbb{A}_{\mathbb{C}}^1 \setminus \{0\}})^{\mathrm{an}}\right) \xrightarrow{\substack{\text{quasi-isomorphic}\\ \sim}}  L_m^{\lambda}.
$$
Since the direct system (in category of coherent (finitely generated) $\D_{\mathbb{A}_{\mathbb{C}}^1}$-modules)
\begin{equation*}
    \begin{array}{ccccccc}
        K_{m-1}^{\alpha} & \lhook\joinrel\xrightarrow{\quad} & K_m^{\alpha}                              & \lhook\joinrel\xrightarrow{\quad} & K_{m+1}^{\alpha}                            & \lhook\joinrel\xrightarrow{\quad} & \cdots \\[1.5ex]
        \overline P      & \xmapsto{\quad}                   & \overline{P \cdot (t\partial_t - \alpha)} & \xmapsto{\quad}                   & \overline{P \cdot (t\partial_t - \alpha)^2} & \xmapsto{\quad}                   & \cdots
    \end{array}
\end{equation*}
induces a direct system of local systems
(categorical equivalence is given by functors, so it also transports morphisms),
arrows should be viewed as $\pi_1((\mathbb{A}_{\mathbb{C}}^1 \setminus \{0\})^\mathrm{an}, p)$-\textbf{equivariant} $\C$-linear maps by Prop.~\ref{prop:classification-of-local-system}:
\begin{equation*}
    \begin{array}{ccccccc}
        L_{m-1}^{\lambda}=\bigoplus \limits_{i=0}^{m-2}\C\cdot v_i & \lhook\joinrel\xrightarrow{\quad} & L_m^{\lambda}=\bigoplus \limits_{i=0}^{m-1}\C\cdot v_i & \lhook\joinrel\xrightarrow{\quad} & L_{m+1}^{\lambda}=\bigoplus \limits_{i=0}^{m}\C\cdot v_i & \lhook\joinrel\xrightarrow{\quad} & \cdots \\[1.5ex]
        v_i                                                        & \xmapsto{\quad}                   & v_i                                                    & \xmapsto{\quad}                   & v_i                                                      & \xmapsto{\quad}                   & \cdots
    \end{array}
\end{equation*}
To see why the morphism are $\pi_1((\mathbb{A}_{\mathbb{C}}^1 \setminus \{0\})^\mathrm{an}, p)$-equivariant,
since $J_{\alpha, m}$ is an upper-triangular matrix, then $e^{-2\pi iJ_{\alpha,m}}$ is also upper-triangular, so each $L^\lambda_{i}\subseteq L^\lambda_{i+1}$ is a sub $\pi_1((\mathbb{A}_{\mathbb{C}}^1 \setminus \{0\})^\mathrm{an}, p)$-representation.

\bigskip

Let $X = \text{Spec } R$ be a smooth affine algebraic variety over $\mathbb{C}$, and $f \in R$,
a regular function, we have the following commutative diagram:
\begin{center}
    \begin{tikzcd}
        {(f=0)\atop \text{divisor}} & X && {\mathbb{A}^1_{\mathbb{C}}} \\
        \\
        & {U=X\backslash(f=0)} && {\mathbb{A}^1_{\mathbb{C}}\backslash \{0\}}
        \arrow["i", hook, from=1-1, to=1-2]
        \arrow["f", from=1-2, to=1-4]
        \arrow["j", hook', from=3-2, to=1-2]
        \arrow["f", from=3-2, to=3-4]
        \arrow[hook', from=3-4, to=1-4]
    \end{tikzcd}
\end{center}
where the morphism $f$ is given by the ring map:
\begin{align*}
    \C[t] & \xrightarrow{\quad} R                       \\
    \C    & \xrightarrow{\mathrm{id}_\C} \C \subseteq R \\
    t     & \xmapsto{\quad} f
\end{align*}
And also note that the locus of $f$ (as an \textbf{analytic set} in complex manifold $X^\mathrm{an}$) can be \textbf{non-smooth}.

We define
$$
    \Psi_{f,\lambda} = \lim_{\substack{\longrightarrow \\ m \to \infty}} i^{-1} Rj_\ast L_m^{\lambda}
$$
which is called the $\lambda\textbf{-nearby cycles}$ of $f$.
And $T$-action ($\pi_1((\mathbb{A}_{\mathbb{C}}^1 \setminus \{0\})^\mathrm{an})$-action) on $L_m^{\lambda}$ induces $T$-action on $\Psi_{f,\lambda}$.



\begin{remark}
    Firstly, note that since $\{f=0\}$ may not smooth, so $i^{-1}$ is the sheaf-theoretic pull-back between topological (ringed) spaces.
    Which is somewhat different from scheme case.

    Secondly, $j_{\ast}$ is the direct image functor,
    and it's a left exact functor,
    we write $Rj_\ast$ the right derived functor of $j_{\ast}$,
    which can be viewed as the complex of sheaves obtained by taking the \textbf{injective resolution} of $L_{m}^{\lambda}$ and applying $j_{\ast}$.
    Since $j_\ast$ maybe have nontrivial cohomology even in a non-$\D$-module-theoretic context, we choose to use a derived functor.
\end{remark}

The above definition is a topological definition, which is a bit differed from Deligne's. And total cycles is defined by combining all information from nearby cycles.


On the other hand, recall Example~\ref{Example 8.11}, we have
$$
    \mathcal{D}_{X}[s]f^{s} \subseteq R\left[\frac{1}{f}\right][s] \cdot f^{s}.
$$

For closed point $\alpha \in \mathbb{A}_{\mathbb{C}}^{1} \simeq \operatorname{Spec} \mathbb{C}[s]$, we use $\mathfrak{m}_{\alpha}$ to denote its maximal ideal. We have proved in Example \ref{Example 8.11}:

\begin{remark}
    We now justify the symbol $R_f[s]f^s$,
    since in Rmk.~\ref{remark:second_treatment_of_example_8.11} we only get the module $R_f[\partial_t]$,
    and $f^s$ represent the canonical section $1\otimes 1$.

    One should note that the notation $R_f[\partial_t]$ (or more precisely $\C[\partial_t]\otimes_\C R_f$) itself doesn't imply any module structure.
    If we want to make it an $R_f$ (left or right) module, we should induce the action from its left $\D_Y$-module structure.

    The Step 5 Block A of Rmk.~\ref{remark:second_treatment_of_example_8.11} gives us an answer.
    We have a projection
    \[
        p: Y = X\times_{\Spec \C}\mathbb{A}^1_{\C} \to X
    \]
    which makes $\Ox$ a subring of $\Oy$.
    So we get a left $R$ module structure by
    \[
        R=\Ox\xhookrightarrow{p^\sharp} \Oy \subseteq \D_Y.
    \]

    Since before we pushforward left $\D_X$-module $R_f$ along $\Gamma_f$,
    $f\in\D_X$ already acts bijectively (left and right) on $R_f$,
    which makes $t\in\D_Y$ act bijectively on $\Gamma_{f,+}(R_f) = \C[\partial_t] \otimes_{\C} R_f$ \cite[Lemma~2.1.5]{PopaDMBG2021}.
    Then by Step 5 Block A+D of Rmk.~\ref{remark:second_treatment_of_example_8.11},
    we discover $f\in R$ also act bijectively on $\Gamma_{f,+}(R_f)$.
    So we finally give a left $R_f$ module structure on $\C[\partial_t] \otimes_{\C} R_f$.

    Then we discover $s = -t\partial_t$ and $R$ commute in $\D_Y$, this is because we have
    \[
        [\partial_t, g] = \partial_t(g) = 0 \quad (\text{see  \cite[Prop~2.25(iii)]{mustata2023dmodules}})
    \]
    By universal property of polynomial ring, we have $R_f[s]$ embedded in $\D_Y$ (as a $\C$-subalgebra),
    if we map formal symbol $s$ to $-t\partial_t \in \D_Y$.

    \medskip

    Then we can act $R_f[s]$ onto $1\otimes 1 \in \Gamma_{f,+}(R_f)$, then we claim:
    \begin{enumerate}
        \item this action is faithful, so we get a free rank 1 $R_f[s]$ submodule in $\Gamma_{f,+}(R_f)$.
        \item This submodule is the whole of $\Gamma_{f,+}(R_f)$.
    \end{enumerate}
    We denote it as $R_f[s]f^s$.
    Then we can give the abstractly constructed $R_f[s]$ a left $\D_Y$ module structure by pulling back the one on $\Gamma_{f,+}(R_f)$.
\end{remark}

\bigskip

\begin{remark}\label{rmk:base_change_of_R_f[s]_f^s}
    We then try to define $R_f(s)f^s$ and $R_f[s]_{\frak{m}_\alpha}f^s$, the construction should satisfy:
    \begin{enumerate}
        \item From (the proof in) the corollary~\ref{cor:bernstein_beilinson} below, which suggests that this construction should be functorial.
        \item Be consistent with \cite[The $\D_{U_s}$ module $\mathcal{M}$ defined in Proof of Theorem 2.2.1.]{PopaDMBG2021}.
    \end{enumerate}

    This is equivalent to say that we have a ring ($\C$-algebra) map
    \[
        \D_Y \to \End_\C(R_f[s])
    \]
    And we have a composition
    \[
        \D_X\otimes_\C \C[s] \to \D_Y \to \End_\C(R_f[s])
    \]
    by assigning $s$ to $-t\partial_t\in\D_Y$

    \medskip

    We define $R_f(s) := R_f[s]\otimes_{\C[s]}\C(s)$ and $R_f[s]_{\frak{m}_\alpha} := R_f[s]\otimes_{\C[s]}\C[s]_{\frak{m}_\alpha}$.
    Before giving it a left $\D$ module structure, we discover that this construction is functorial (in the category of $\C[s]$-modules).

    \medskip

    Then we give it a left $\D$ module structure, this structure should also be functorial.
    \begin{enumerate}
        \item $R_f(s)$ should have a left $\D_{R(s)/\C(s)} = \D_{R/\C}\otimes_\C \C[s]\otimes_{\C[s]}\C(s)$ module structure, see Lem.~\ref{lem:base_change_of_D_modules}.
        \item $R_f[s]_{\frak{m}_\alpha}$ should have a left $\D_{R[s]_{\frak{m}_\alpha}/\C[s]_{\frak{m}_\alpha}} = \D_{R/\C}[s]\otimes_{\C[s]} \C[s]_{\frak{m}_\alpha}$ module structure, but one should note $\C[s]_{\frak{m}_\alpha}$ is not a field.
    \end{enumerate}
    Thanks to the construction in the proof in Lem.~\ref{lem:base_change_of_D_modules}, the base change of the \textbf{ring of \underline{relative} differential operator} is functorial.

    \medskip

    To give a (left) module structure is equivalent to give the ring map
    \begin{align*}
        \D_{R/\C}[s]\otimes_{\C[s]}\C(s) & \to \End_\C(R_f[s]\otimes_{\C[s]}\C(s))                    \\
        P \otimes a                      & \mapsto \left( u \otimes b \mapsto P(u) \otimes ab \right)
    \end{align*}
    and
    \begin{align*}
        \D_{R/\C}[s]\otimes_{\C[s]}\C[s]_{\frak{m}_\alpha} & \to \End_\C(R_f[s]\otimes_{\C[s]}\C[s]_{\frak{m}_\alpha})  \\
        P \otimes a                                        & \mapsto \left( u \otimes b \mapsto P(u) \otimes ab \right)
    \end{align*}
    From the two formulas above, one see this definition of ring action (on module) is functorial. Where $u\mapsto P(u)$ is just
    \[
        \D_{R/\C} \otimes_\C \C[s] = \D_X\otimes_\C \C[s] \to \D_Y \to \End_\C(R_f[s])
    \]
    mentioned above.
\end{remark}

\begin{theorem}
    There exists $b(s) \in \mathbb{C}[s]$ s.t.
    \[
        b(s) \quad \mathrm{ kills } \quad \frac{\mathcal{D}_{X}[s]f^{s}}{\mathcal{D}_{X}[s]f^{s+1}}
    \]
    see Thm.~\ref{Theorem 8.4} or Ex.~\ref{Example 8.11} or Rmk.~\ref{remark:second_treatment_of_example_8.11}.
\end{theorem}



\begin{corollary}[Bernstein-Beilinson]\label{cor:bernstein_beilinson}
    Let $\mathbb{C}(s)$ be the fractional field of $\mathbb{C}[s]$, and $\mathcal{D}_{X}(s) = \mathcal{D}_{X} \otimes_{\mathbb{C}} \mathbb{C}(s)$; for all $\alpha \in \mathbb{A}_{\mathbb{C}}^{1}$ closed point, we obtain that:
    \begin{enumerate}
        \item $\mathcal{D}_{X}(s) f^{s+k} \simeq \mathcal{D}_{X}(s) f^{s}, $ for all $  k \in \mathbb{Z}$.
        \item There exists $k_{0} \geqslant 0$ s.t.
              \begin{itemize}
                  \item $\mathcal{D}_{X}[s]_{\mathfrak{m}_{\alpha}} f^{s-k} = \mathcal{D}_{X}[s]_{\mathfrak{m}_{\alpha}} f^{s-k_{0}}, $ for all $ k > k_{0}$
                  \item $\mathcal{D}_{X}[s]_{\mathfrak{m}_{\alpha}} f^{s+k} = \mathcal{D}_{X}[s]_{\mathfrak{m}_{\alpha}} f^{s+k_{0}},$ for all $  k > k_{0}$
              \end{itemize}
    \end{enumerate}

\end{corollary}

\begin{remark}\label{rmk:D_X(s)f^s=R_f(s)f^s_and_D_X[s]_m_alpha_f^s=R_f[s]_m_alpha_f^s}
    Noted that here  $\mathcal{D}_{X}(s) f^{s} = R[\frac{1}{f}](s) f^{s},\ \mathcal{D}_{X}[s]_{\mathfrak{m}_{\alpha}} f^{s} = R[\frac{1}{f}][s]_{m_{\alpha}} f^{s}$.
\end{remark}

\begin{proof}
    \mbox{}

    $1)$ Since $\mathbb{C}(s)$ can be view as the localization at $(0)$ of $\mathbb{C}[s]$,
    we have $\mathbb{C}(s)$ is flat over $\mathbb{C}[s]$,
    thus
    \[
        \frac{\mathcal{D}_{X}[s]f^{s}}{\mathcal{D}_{X}[s]f^{s+1}}\otimes \mathbb{C}(s)=\frac{\mathcal{D}_{X}(s)f^{s}}{\mathcal{D}_{X}(s)f^{s+1}}
    \]
    Note that
    \begin{enumerate}
        \item To commute tensor product with quotient, the ring map $\mathbb{C}[s]\to \mathbb{C}(s)$ is not needed to be flat, all tensor product is right exact.
        \item To make the incuced $\D_X(s)$-structure the same on LHS and RHS, the $\D_X(s)$-structure should be functorial, which is required in Rmk.~\ref{rmk:base_change_of_R_f[s]_f^s}.
    \end{enumerate}

    From the $\D_X(s)$-structure mentioned in Rmk.~\ref{rmk:base_change_of_R_f[s]_f^s}, we have
    \[
        b(s+k) \quad\mathrm{kills} \quad\frac{\mathcal{D}_{X}(s)f^{s}}{\mathcal{D}_{X}(s)f^{s+1}} \implies \D_{R/\C}[s]\otimes_{\C[s]} \C(s) \ni b(s+k)\otimes 1  \quad\mathrm{kills } \quad\frac{\mathcal{D}_{X}[s]f^{s}}{\mathcal{D}_{X}[s]f^{s+1}} \otimes_{\C[s]} \C(s)
    \]
    And $b(s+k) \neq 0$ is invertible in $\mathbb{C}(s)$, so
    \[
        b(s+k) \quad\mathrm{kills} \quad\frac{\mathcal{D}_{X}(s)f^{s}}{\mathcal{D}_{X}(s)f^{s+1}} \implies \frac{\mathcal{D}_X(s) f^{s+k}}{\mathcal{D}_X(s) f^{s+k-1}} = 0.
    \]

    $2)$  Note that Bernstein-Sato polynomial $b(s)\in\D_X[s]$ has coefficients in $\C$, see Thm.~\ref{Theorem 8.4}.

    Since the number of roots of $b(s)$ is finite,
    there exists $k_0 \ge 0$, $b(\alpha \pm k) \neq 0 \in\mathbb{C}$ for all $k > k_0$,
    then $b(s \pm k) \notin \mathfrak{m}_\alpha$ for all $k > k_0$.

    So $b(s \pm k)$ is invertible in $\mathbb{C}[s]_{\mathfrak{m}_\alpha}$ for all $k > k_0$,
    similarly, we get
    \[
        \frac{\mathcal{D}_X[s]_{\mathfrak{m}_\alpha} f^{s \pm k}}{\mathcal{D}_X[s]_{\mathfrak{m}_\alpha} f^{s \pm k + 1}} = 0.
    \]
\end{proof}

\begin{definition}
    For all $\alpha \in \mathbb{C}$:
    \begin{enumerate}
        \item $\mathcal{D}_X[s]_{\mathfrak{m}_\alpha} f^{s+k},  k \gg 0$ is stabled, called the $\textbf{minimal extension}$ of $R[\frac{1}{f}][s]_{\mathfrak{m}_\alpha} f^s$.
        \item $\Psi_{f,\alpha} := \frac{\mathcal{D}_X[s]_{\mathfrak{m}_\alpha} f^{s-k}}{\mathcal{D}_X[s]_{\mathfrak{m}_\alpha} f^{s+k}}$ is called the $\alpha\textbf{-nearby cycles}$ of $R$ along $f$.

              \cite[Definition~3.3.2.]{PopaDMBG2021} defines that $\alpha$-nearby cycles is a $\mathcal{D}_X$-module, NOT a $\mathcal{D}_X[s]$-module.
        \item $\Lambda := \mathbb{Z} + \{\text{roots of } b_{f}(s)\}$.
    \end{enumerate}
\end{definition}

\begin{remark}
    For one wonder why here we \emph{define} the minimal extension,
    see \cite[Theorem~2.10]{Wu2020ComparisonNearbyCyclesBfunctions}, which use
    \cite[Def~3.4.1]{hotta2008d} to define the minimal extension and claim the equivalence of the two.

    In $\mathcal{D}$-module theory,
    an extension is usually about given a module over an open subset $U$,
    and extending it to the whole space $X$ to get a $\mathcal{D}_{X}$-module,
    and we shall pay attention on how it extends around $f=0$.
    Take $U=\mathrm{Spec}\, R_{f}$,
    if the module over $U=X\backslash(f=0)$ can be viewed as $R$-module,
    then it's an extension,
    and we call $R_{f}[s]_{\mathfrak{m}_{\alpha}}f^{s}$  $\textbf{maximal extension}$.
\end{remark}

\begin{proposition}
    \quad

    \begin{enumerate}
        \item $s - \alpha$ acts on $\Psi_{f,\alpha}$ nilpotently
        \item $\Psi_{f,\alpha} \cong \Psi_{f,\alpha+1}$ (Note: a $\D_X$-module homomorphism, not a $\D_X[s]$-module homomorphism)
        \item $\Psi_{f,\alpha} \neq 0$ if and only if $\alpha \in \Lambda$
        \item $\Psi_{f,\alpha}$ is a regular holonomic $D_X$-module.
    \end{enumerate}
\end{proposition}
\begin{proof}
    \mbox{}

    \begin{enumerate}
        \item The descending chain of submodules
              \[
                  \mathcal{D}_{X}[s]_{\mathfrak{m}_{\alpha}}f^{s+k_{0}}\subseteq\mathcal{D}_{X}[s]_{\mathfrak{m}_{\alpha}}f^{s+k_{0}-1}\subseteq\dots\subseteq \mathcal{D}_{X}[s]_{\mathfrak{m}_{\alpha}}f^{s-k_{0}}
              \]
              induces a finite filtration on the quotient module
              \[
                  \frac{\mathcal{D}_{X}[s]_{\mathfrak{m}_{\alpha}}f^{s-k_{0}}}{\mathcal{D}_{X}[s]_{\mathfrak{m}_{\alpha}}f^{s+k_{0}}},
              \]
              whose successive quotients are of the form
              \[
                  \frac{\mathcal{D}_{X}[s]_{\mathfrak{m}_{\alpha}}f^{s+i}}{\mathcal{D}_{X}[s]_{\mathfrak{m}_{\alpha}}f^{s+i+1}}.
              \]
              For each such quotient, there exists a polynomial $b_{i}(s)$ that annihilates it (before taking localization). Because $s-\alpha$ is the only non-invertible irreducible polynomial in the local ring $\mathbb{C}[s]_{\mathfrak{m}_{\alpha}}$, each $b_{i}(s)$ takes the form $u_{i}(s-\alpha)^{n_{i}}$ for some unit $u_{i} \in \mathbb{C}[s]_{\mathfrak{m}_{\alpha}}$.
              By taking $n = \sum n_{i}$, the element $(s-\alpha)^{n}$ annihilates the entire module $\frac{\mathcal{D}_{X}[s]_{\mathfrak{m}_{\alpha}}f^{s-k_{0}}}{\mathcal{D}_{X}[s]_{\mathfrak{m}_{\alpha}}f^{s+k_{0}}}$. Hence, $s-\alpha$ acts nilpotently on it.

        \item Recall that $t$ acts on $R_f[s]f^s$ bijectively. Consider the $\mathcal{D}_{X}$-linear automorphism of $\mathcal{D}_{X}(s)f^{s}$ defined by the shift
              \[
                  T_{-1}\big(P(s)f^{s}\big) = P(s-1)f^{s-1}.
              \]
              For any element $m$, we have
              \[
                  T_{-1}(s \cdot m) = (s-1) \cdot T_{-1}(m), \quad\text{and generally}\quad T_{-1}(P(s) \cdot m) = P(s-1) \cdot T_{-1}(m).
              \]
              This means $T_{-1}$ is not $\D_X[s]$-linear.

              When localizing at $\mathfrak{m}_{\alpha}$,
              a polynomial $g(s)$ becomes invertible if $g(\alpha) \neq 0$.
              Under the shift, the action of $g(s)$ becomes the action of $g(s-1)$, evaluated at $s = \alpha+1$ is also non-zero, making it invertible in $\mathbb{C}[s]_{\mathfrak{m}_{\alpha+1}}$.
              So we can safely define the map below:
              \begin{align*}
                  T_{-1}: \mathcal{D}_{X}[s]_{\mathfrak{m}_{\alpha}}f^{s-k} & \stackrel{\sim}{\longrightarrow} \mathcal{D}_{X}[s]_{\mathfrak{m}_{\alpha+1}}f^{s-k-1} \\
                  \frac{m}{g(s)}                                            & \longmapsto \frac{T_{-1}(m)}{g(s-1)}
              \end{align*}
              for any integer $k$. And one should verify that it is indeed an isomorphism.

              For $k \gg 0$, the sequence of modules stabilizes. Passing to the respective quotients, this induces the desired $\mathcal{D}_X$-module isomorphism $\Psi_{f,\alpha} \xrightarrow{\sim} \Psi_{f,\alpha+1}$.

        \item Since $\alpha \not\in\Lambda$, $\alpha$ is not a root of $b(s+k)$, then $b(s+k)$ is a unit for all $k$, thus $\Psi_{f,\alpha}=0$.

              If $\Psi_{f,\alpha}\neq 0$,
              there exists $i$ such that $\frac{\mathcal{D}_{X}[s]_{\mathfrak{m}_{\alpha}}f^{s+i}}{\mathcal{D}_{X}[s]_{\mathfrak{m}_{\alpha}}f^{s+i+1}}\neq 0$,
              and then $\alpha$ is a root of $b(s+i)$ such that $b(s+i)$ is not invertible in $\mathbb{C}[s]_{\mathfrak{m}_{\alpha}}$. Thus $\alpha\in\Lambda$.

        \item $M_{\alpha} = \mathbb{C}[t, \frac{1}{t}] \cdot t^{\alpha}$ (see Ex.~\ref{ex:first-examples-of-d-modules}) is regular
              (at least analytically at the stalk of $0$, see Rmk.~\ref{rmk:indecomposable_objects_in_RH_D} and Prop.~\ref{prop:classification-unipotent-indecomposable-regular-holonomic},
              to be verified algebraically it should be embedded into $\mathbb{P}^1_\C$, see Rmk.~\ref{rmk:algebraic_KK_regularity_holomonic})
              holonomic $\mathcal{D}_{\mathbb{A}^{1}_{\mathbb{C}}}$-module.

              Consider
              \begin{align*}
                  f: X   & \longrightarrow \mathbb{A}^{1}_{\mathbb{C}} \\
                  \alpha & \longmapsto t=f(\alpha)
              \end{align*}
              which is induced by
              \begin{align*}
                  \mathbb{C}[t] & \longrightarrow R \\
                  t             & \longmapsto f
              \end{align*}
              We admit the fact that the $\D$-module theoretical (derived) \textbf{inverse image} functor \cite[Sect.~1.5, It is the derived one that comptible with Thm.~6.1.5]{hotta2008d} preserves (algebraic we use here) regular holonomicity \cite[Thm.~6.1.5]{hotta2008d}, then
              \[
                  f^{*} M_{\alpha}:=M_{\alpha}\otimes_{\mathbb{C}[t]} R\cong R_{f}\cdot f^\alpha= R[\frac{1}{f}] f^{\alpha}
              \]
              makes $R_{f}\cdot f^\alpha$ a regular holonomic $\mathcal{D}_{X}$-module.

              With $(2)$, we can assume $\alpha\neq 0$, otherwise, we just shifts $\alpha$ to $\alpha+1$.
              We have a short exact sequence:
              \[
                  0 \longrightarrow \frac{R_f[s]_{\mathfrak{m}_{\alpha}} f^{s}}{(s-\alpha) R_f[s]_{\mathfrak{m}_{\alpha}} f^{s}} \xrightarrow{\;(s-\alpha)\cdot\;} \frac{R_f[s]_{\mathfrak{m}_{\alpha}} f^{s}}{(s-\alpha)^2 R_f[s]_{\mathfrak{m}_{\alpha}} f^{s}} \xrightarrow{\;\text{proj}\;} \frac{R_f[s]_{\mathfrak{m}_{\alpha}} f^{s}}{(s-\alpha) R_f[s]_{\mathfrak{m}_{\alpha}} f^{s}} \longrightarrow 0
              \]
              where the injective mapping is given by \textbf{left} multiplication by $(s-\alpha)$ (i.e., sending the residue class $\overline{m}$ to $\overline{(s-\alpha)m}$), and the surjective mapping is the natural canonical projection.

              Here, the isomorphism
              \[
                  \frac{R_f[s]_{\mathfrak{m}_{\alpha}} f^{s}}{(s-\alpha) R_f[s]_{\mathfrak{m}_{\alpha}} f^{s}} \cong R_{f} \cdot f^{\alpha}
              \]
              arises because quotienting out the $(s-\alpha)$-multiples algebraically corresponds to taking the evaluation map $s \mapsto \alpha$.
              \begin{enumerate}
                  \item The ring $\mathbb{C}[s]_{\mathfrak{m}_{\alpha}}$ is a local ring with (its unique) maximal ideal $(s-\alpha)$ reduces canonically to $\mathbb{C}$, collapsing the coefficients from $R_f[s]_{\mathfrak{m}_{\alpha}}$ down to $R_{f}$.
                  \item And the formal symbol $f^{s}$ becomes $f^{\alpha}$, this step didn't actually do evaulation, since they are \textbf{both} the canonical section $1\otimes 1$ in their respective presentations. So the real thing goes $1\otimes 1 \mapsto 1\otimes 1$.
              \end{enumerate}
              Thus, this quotient module is precisely the regular holonomic $\mathcal{D}_X$-module $R_{f} \cdot f^{\alpha}$.

              Since regular holonomic $\mathcal{D}_{X}$-modules form an abelian category \cite[Combine Prop.~3.1.2 and Def.~6.1.1]{hotta2008d}, we have $\frac{R_{f}[s]_{\mathfrak{m}_{\alpha}}f^{s}}{(s-\alpha)^{2}R_{f}[s]_{\mathfrak{m}_{\alpha}}f^{s}}$ is regular holonomic.

              And further, we obtain the generalized short exact sequence:
              \[
                  0 \longrightarrow \frac{R_f[s]_{\mathfrak{m}_{\alpha}} f^{s}}{(s-\alpha) R_f[s]_{\mathfrak{m}_{\alpha}} f^{s}} \xrightarrow{\;(s-\alpha)^N\cdot\;} \frac{R_f[s]_{\mathfrak{m}_{\alpha}} f^{s}}{(s-\alpha)^{N+1} R_f[s]_{\mathfrak{m}_{\alpha}} f^{s}} \xrightarrow{\;\text{proj}\;} \frac{R_f[s]_{\mathfrak{m}_{\alpha}} f^{s}}{(s-\alpha)^N R_f[s]_{\mathfrak{m}_{\alpha}} f^{s}} \longrightarrow 0
              \]
              By induction, $\frac{R_f[s]_{\mathfrak{m}_{\alpha}} f^{s}}{(s-\alpha)^{N} R_f[s]_{\mathfrak{m}_{\alpha}} f^{s}}$ is regular holonomic for all $N$.

              By $(1)$, the action of $(s-\alpha)$ on $\Psi_{f,\alpha} = \frac{\mathcal{D}_{X}[s]_{\mathfrak{m}_{\alpha}}f^{s-k}}{\mathcal{D}_{X}[s]_{\mathfrak{m}_{\alpha}}f^{s+k}}$ is nilpotent, meaning there exists some large enough integer $N \gg 0$ such that $(s-\alpha)^N$ entirely annihilates $\Psi_{f,\alpha}$.
              Because of this annihilation, the natural surjective $\mathcal{D}_X[s]_{\mathfrak{m}_{\alpha}}$-module morphism
              \[
                  \mathcal{D}_{X}[s]_{\mathfrak{m}_{\alpha}}f^{s-k} \xrightarrow{\;\text{proj}\;} \frac{\mathcal{D}_{X}[s]_{\mathfrak{m}_{\alpha}}f^{s-k}}{\mathcal{D}_{X}[s]_{\mathfrak{m}_{\alpha}}f^{s+k}} = \Psi_{f,\alpha}
              \]
              factors through the quotient by the submodule generated by $(s-\alpha)^N$. This explicitly means that the map descends to a well-defined surjective homomorphism:
              \[
                  \frac{\mathcal{D}_X[s]_{\mathfrak{m}_{\alpha}} f^{s-k}}{(s-\alpha)^N \mathcal{D}_X[s]_{\mathfrak{m}_{\alpha}} f^{s-k}} \twoheadrightarrow \Psi_{f,\alpha}.
              \]
              Since one can verify
              \begin{enumerate}
                  \item $\mathcal{D}_{X}[s]_{\mathfrak{m}_{\alpha}}f^{s-k} \cong \mathcal{D}_{X}[s]_{\mathfrak{m}_{\alpha}}f^{s}$,
                  \item Rmk.~\ref{rmk:D_X(s)f^s=R_f(s)f^s_and_D_X[s]_m_alpha_f^s=R_f[s]_m_alpha_f^s},
                        i.e.,
                        $\mathcal{D}_X[s]_{\mathfrak{m}_{\alpha}} f^{s}$ as a submodule $R_f[s]_{\mathfrak{m}_{\alpha}} f^{s}$ actually equals it.
              \end{enumerate}
              Thus, we can view $\Psi_{f,\alpha}$ as a quotient of the $\mathcal{D}_{X}$-module $\frac{R_f[s]_{\mathfrak{m}_{\alpha}} f^{s}}{(s-\alpha)^N R_f[s]_{\mathfrak{m}_{\alpha}} f^{s}}$.
              So $\Psi_{f,\alpha}$ itself must also be a regular holonomic $\mathcal{D}_{X}$-module.
    \end{enumerate}

\end{proof}


\section{Comparison theorem}
Our next goal is to prove the following theorem:
\begin{theorem}[Comparison] \label{Theorem 9.5}
    For all $\alpha \in \mathbb{C}$, $\lambda = e^{2\pi i\alpha}$, $$\mathrm{DR}(\Psi_{f,-\alpha}) \simeq i_* \Psi_{f,\lambda}[-1]$$such that $e^{2\pi is} = \lambda^{-1}\cdot e^{2\pi i(s+\alpha)}$ is the corresponding $T$-action.
\end{theorem}
\begin{remark}
    with $(s-\alpha)$ acts on $\Psi_{f,\alpha}$ nilpotent, $e^{2\pi i(s+\alpha)}$ will be a well-defined operator for $\Psi_{f,-\alpha}$.
\end{remark}


Define $N_m^{\alpha,k} = \bigoplus_{\ell=1}^m \mathcal{D}_X[s] f^{s-k} \cdot e_\ell^{\alpha}$, know that $s=-t\partial_{t}$,  $s \cdot (\eta \cdot e_\ell^{\alpha}) = (s\eta) \cdot e_\ell^{\alpha}+\eta(-\alpha e^{\alpha}_{\ell}-e^{\alpha}_{\ell-1})=(s-\alpha)\eta \cdot e^{\alpha}_{\ell}-\eta \cdot e^{\alpha}_{\ell-1}$, so $N_m^{\alpha,k}$ is also an $\mathcal{D}_X[s]$-mod.

We define $b_f(s)$ to be the monic polynomial of least degree among $b(s)$ such that $$b(s) \cdot \frac{\mathcal{D}_X[s] f^s}{\mathcal{D}_X[s] f^{s+1}} = 0$$called the $\textbf{b-function} \text{ or } \textbf{Bernstein-Sato polynomial of }f$.

\begin{lemma}\label{lemma 9.6}
    \quad

    \begin{enumerate}
        \item $\mathrm{DR}\left(R_f[s] f^s \otimes^L_{\mathbb{C}(s)} \mathbb{C}_\alpha\right) \simeq \mathrm{DR}\left(R_f[s] f^{s+\alpha} \otimes^L_{\mathbb{C}(s)} \mathbb{C}_\alpha\right)\simeq R j_* (f^{-1} L_1^\lambda)$
        \item $\mathrm{DR}\left(\mathcal{D}_X[s] f^{s+k} \otimes^{L}_{\mathbb{C}(s)} \mathbb{C}_\alpha\right) \simeq \mathrm{DR}\left(\mathcal{D}_X[s] f^{s+k+\alpha} \otimes^{L}_{\mathbb{C}(s)} \mathbb{C}_\alpha\right)\simeq j_! (f^{-1} L_1^\lambda), \quad k \gg |\alpha|$
    \end{enumerate}
    where $\mathbb{C}_\alpha$ is the residue field of $\alpha \in \mathbb{A}_{\mathbb{C}}^1$, and $\lambda = e^{2\pi i \alpha}$.
\end{lemma}
\begin{proof}
    Take $0\rightarrow \mathbb{C}[s]\xrightarrow{s-\alpha} \mathbb{C}[\alpha]\rightarrow \frac{\mathbb{C}[s]}{s-\alpha}=\mathbb{C}_{\alpha}$ free resolution, then we obtain that $R_f[s]f^s \otimes^L_{\mathbb{C}(s)} \mathbb{C}_\alpha$ $$\stackrel{\text{q.i}}{\simeq}[ R_f[s]f^s \xrightarrow{s-\alpha} R_f[s]f^s ]\stackrel{\text{q.i}}{\simeq}[R_{f}[s]f^{s+\alpha}\xrightarrow{s} R_f[s]f^{s+\alpha}]\simeq R_f f^\alpha.$$with $U\xrightarrow{f}\mathbb{A}^{1}_{\mathbb{C}}\backslash\{0\}$, $\mathrm{DR}(R_f \cdot f^\alpha)|_U \simeq f^{-1} L_1^\lambda$, using log resolution of $(X,f=0)$, one can prove $\mathrm{DR}(R_f \cdot f^\alpha) \simeq R j_*(f^{-1} L_1^\lambda)$.
    Using resolution of singularities and duality, one can prove 2.

\end{proof}
\begin{lemma} \label{lemma 9.7}
    $$(b_f(s-k-\alpha))^m \text{ kills } N_m^{\alpha,k}/N_m^{\alpha,k-1}$$
\end{lemma}
\begin{proof}
    By construction, we have a SES: $$0\longrightarrow N_{m-1}^{\alpha,k} \longrightarrow N_m^{\alpha,k} \longrightarrow Q \longrightarrow 0$$with $Q=\mathcal{D}_{X}[s]f^{s-k}\bar{e}_{m-1}^{\alpha}$, with $s\cdot \bar{e}^\alpha_{m-1}=(s-\alpha)\overline{e}^\alpha_{m-1}$, we choose $s-\alpha$ to replace $s$, then  $Q\simeq \mathcal{D}_{X}[s]f^{s-k-\alpha}$.

    Since $b_{f}(s-k-\alpha)$ kills $\mathcal{D}_{X}[s]f^{s-k-\alpha}$, and by induction, we have $(b_f(s-k-\alpha))^{m -1}$ kill $N_{m-1}^{\alpha,k}$, done.
\end{proof}
\begin{proposition}
    There exists $k_0 > 0$ such that for all $k \geq k_0$, we obtain that:
    \begin{enumerate}
        \item $\mathrm{DR}(N_m^{-\alpha,k} \xrightarrow{s} N_{m}^{-\alpha,k}) \simeq R j_*\left(f^{-1} L_m^{\lambda}\right)$;
        \item $\mathrm{DR}(N_m^{-\alpha,-k} \xrightarrow{s} N_{m}^{-\alpha,-k})\simeq j_!\left(f^{-1} L_m^{\lambda}\right)$.
    \end{enumerate}
\end{proposition}
\begin{proof}
    Let $\underline{\mathfrak{m}}_0$ be the maximal ideal of $0 \in \mathbb{A}_{\mathbb{C}}^1$, then $$[N_{m}^{-\alpha,k} \xrightarrow{s} N_{m}^{-\alpha,k}]\stackrel{\text{q.i}}{\simeq}N^{-\alpha,k} \otimes^{L}_{\mathbb{C}[s]} \mathbb{C}_{0} \stackrel{\text{q.i}}{\simeq} [N^{-\alpha,k}_{m,\underline{\mathfrak{m}}_{0}} \xrightarrow{s}N^{-\alpha,k}_{m,\underline{\mathfrak{m}}_{0}}]\text{ localization}\atop \stackrel{\text{q.i}}{\simeq}[\bigoplus \limits_{\ell=1}^{m-1} R_{f}[s]f^{s} \cdot e_{\ell}^{-\alpha} \xrightarrow{s} \bigoplus \limits_{\ell=1}^{m-1}R_{f}[s]f^{s} \cdot e_{\ell}^{-\alpha}]\text{ by Lemma \ref{lemma 9.7}}.$$Using Lemma \ref{lemma 9.6} (1) for $\mathrm{DR}(R_{f}[s]f^{s+\alpha}\otimes_{\mathbb{C}[s]}\mathbb{C}_{0})$, by induction, Part (1) is proved. And part (2) is similar.
\end{proof}

\begin{lemma}
    Let $W$ be a $\mathbb{C}$-vector space, and let $\varphi$ be a $\mathbb{C}$-linear operator on $W$ admitting a minimal polynomial.
    Set $W_{\infty} = \bigoplus_{k=0}^{\infty} W \otimes e_k$ and define $$\varphi_{\infty}(w \otimes e_i) = (\varphi - \alpha) \cdot w \otimes e_i - w \otimes e_{i-1} \quad \text{for all }  w \in W, \; (e_{-1} = 0).$$Then $\varphi_{\infty}$ is surjective and $\ker(\varphi_{\infty}) = W_{\alpha}$ is the generalized $\alpha$-eigenspace.
\end{lemma}
\begin{proof}
    Define a map
    \begin{equation}
        W_{\alpha} \longrightarrow \ker(\varphi_{\infty})\atop w \longmapsto \sum\limits_{i \geqslant 0} (\varphi - \alpha)^i w \otimes e_{i}.
        \tag{$\ast$}
        \label{eq}
    \end{equation}
    This map is an isomorphism.

    Then we prove surjectivity: If $w\in W_{\alpha}$, then $\varphi_{\infty}( -\sum\limits_{i>j} (\varphi-\alpha)^{i-j-1}w\otimes e_{i})=w\otimes e_{j}$; If $w\in W_{\alpha}^{\perp}$, then $\varphi_{\infty}( \sum\limits_{i=1}^{j}(\varphi-\alpha)^{-i}w\otimes e_{j-i+1} )=w\otimes e_{j}$.
\end{proof}
\begin{corollary}
    For $k \gg 0$, $$
        \Psi_{f,\alpha}[1] \stackrel{\text{q.i}}{\simeq}\varinjlim_m \left( \frac{N_m^{\alpha,k}}{N_m^{\alpha,-k}} \xrightarrow{s} \frac{N_m^{\alpha,k}}{N_m^{\alpha,-k}} \right).
    $$
\end{corollary}

\begin{proof}[Proof of  (\ref{Theorem 9.5})]
    $$
        \mathrm{DR}(\Psi_{f,-\alpha}[1]) \cong \mathrm{DR}\left( \varinjlim_m  \left( \frac{N_m^{-\alpha,k}}{N_m^{-\alpha,-k}} \xrightarrow{s} \frac{N_m^{-\alpha,k}}{N_m^{-\alpha,-k}} \right) \right)$$
    $$ \cong \varinjlim_m  \mathrm{DR}\left( \frac{N_m^{-\alpha,k}}{N_m^{-\alpha,-k}} \xrightarrow{s} \frac{N_m^{-\alpha,k}}{N_m^{-\alpha,-k}} \right)$$
    $$\simeq \varinjlim_m\ \text{cone}\left( DR\left(N_{m}^{-\alpha,-k} \xrightarrow{s} N_m^{-\alpha,-k}\right) \rightarrow DR\left(N_m^{-\alpha,k} \xrightarrow{s} N_m^{-\alpha,k}\right)\right)$$
    $$\simeq\varinjlim_m\ \text{cone}\left(j_! f^{-1} L_m^{\lambda} \rightarrow R j_* f^{-1} L_m^{\lambda}\right)\simeq \Psi_{f,\lambda}$$

    By $\eqref{eq}$, we consider $$   \Psi_{f,-\alpha} \longrightarrow \ker(-s_{\infty})\atop \eta \longmapsto \sum\limits_{i \geqslant 0} (-s-\alpha)^i \eta\otimes  e_{i}.$$

    Then $s+\alpha = -J_{0,\infty}$, and $e^{2\pi i(s+\alpha)} = e^{-2\pi i J_{0,\infty}}$. Furthermore, we obtain that $$e^{-2\pi i\alpha} \cdot e^{2\pi i(s+\alpha)} = e^{-2\pi i J_{\alpha,\infty}} = T$$
\end{proof}
